% ===== Block C: PDF p.10--15 | 原文 §2.2 Stratification and tautological maps | 唯一编号公式 2.17--2.35 =====
% [原文 p.10]

关于 Deligne--Mumford 模空间的主要结果是：它给出黎曼曲面模空间的一个紧化。

\begin{theorem}[2.4]
  设 $2g-2+n>0$。模空间 $\Mc_{g,n}$
  \begin{itemize}
    \item 是连通的、光滑的、复的、紧的轨形（orbifold），其维数为 $\dim(\Mc_{g,n})=3g-3+n$；
    \item 包含 $\cM_{g,n}$ 作为一个开的稠密子集。
  \end{itemize}
\end{theorem}

% [Tutor Note: 维数公式 $\dim=3g-3+n$ 的两次自检：$g=0,n=3$ 给 0（$\cM_{0,3}$ 是一点）；$g=1,n=1$ 给 1（$j$-线）。另一自检是它与 Teichmüller 空间维数的一半相符：$T_{g,n}$ 的实维数为 $6g-6+2n$，故复维数为 $3g-3+n$。原文只说“轨形”而不说“流形”，是因为稳定曲面可以有非平凡的自同构群（定义 2.3 末条），这些点上的局部模型是 $\C^{N}/\text{有限群}$，对应于把 “orbifold” 译为“轨形”而非“轨道流形”的用法。]

\emph{集合 $\partial\Mc_{g,n}=\Mc_{g,n}\setminus\cM_{g,n}$ 称为模空间的边界。}

既然有了紧空间，我们就可以放心地讨论积分了。更一般地，我们有一个很好的（上）同调代数
\begin{equation}
  (H_{\bullet}(\Mc_{g,n},\Q),\frown)
  \qquad \text{与} \qquad
  (\Ho(\Mc_{g,n},\Q),\smile),
  \tag{2.17}
\end{equation}
其中代数结构分别相对于 cap/cup 积（依次对应于子簇的相交与微分形式的外积）。$\Q$ 系数来自轨形结构；若愿意，也可以放心地取 $\C$ 系数。二者通过 Poincaré 对偶互为对偶：
\begin{equation}
  H^k(\Mc_{g,n},\Q)\cong H_{2(3g-3+n)-k}(\Mc_{g,n},\Q).
  \tag{2.18}
\end{equation}
最重要的是，我们有一个定义良好的基本类，把上同调类对它积分就得到一个数：
\begin{equation}
  \int_{\Mc_{g,n}}\alpha\in\Q,\qquad
  \alpha\in H^{2(3g-3+n)}(\Mc_{g,n},\Q).
  \tag{2.19}
\end{equation}
由于在上同调中取 cap 积（即微分形式的外积）对应于在同调中取 cup 积（即子簇的相交），紧模空间上的积分理论通常被称为\emph{相交理论}。

% [Tutor Note: (2.17) 把 $\smile$ 配给上同调、把 $\frown$ 配给同调，但严格说来 $\frown$ 不是同调上的“代数结构”：cap 积是配对 $H^{p}\otimes H_{q}\to H_{q-p}$，使 $H_{\bullet}$ 成为 $H^{\bullet}$-模，而不是环；只有 cup 积才使 $H^{\bullet}$ 成为环。原文用 “(co)homology algebra” 一词略松，可读作“环与模”。核对方法：$H_{\bullet}$ 上没有自然的结合乘法（相交积只在光滑流形同调上、且需横截性才良定义），而 $H^{\bullet}$ 的 cup 积总是有定义的。]
% [Tutor Note: (2.18) 的移位量 $2(3g-3+n)$ 是 $\Mc_{g,n}$ 的实维数（$=\dim_{\C}\times 2$），所以对偶配对是 $H^{k}\times H^{2\dim_{\C}-k}\to H^{2\dim_{\C}}$。核对：$k=0$ 时得 $H^{0}\cong H_{2(3g-3+n)}$，两者都同构于 $\Q$（紧连通轨形）。]
% [Tutor Note: (2.19) 说明为什么把基本类叫作“基本”：只有最高次 $2(3g-3+n)$ 的类才能与它配对成数，因此只有这些类谈得上“积分”。也正因如此，本讲义后面所有相交数（例如 $\int_{\Mc_{g,n}}\psi_1^{k_1}\cdots\psi_n^{k_n}$）都先要求次数相加等于 $2(3g-3+n)$，否则积分恒为零。]

\subsection{分层结构与重言映射（Stratification and tautological maps）}

在进一步讨论 $\Mc_{g,n}$ 的上同调及其相交理论之前，让我们更仔细地分析这个紧化。需要记住的主要图景如下：$\Mc_{g,n}$ 的大部分点都是落在 $\cM_{g,n}$ 上的光滑黎曼曲面；但通过收缩圈，我们会得到落在边界 $\partial\Mc_{g,n}$ 上的稳定奇异黎曼曲面。做一次这样的操作，就造出一个结点。反复进行这一操作，我们得到的黎曼曲面就越来越奇异。见图 7 的图示。

作为例子，考虑空间 $\Mc_{0,4}$。在边界 $\partial\Mc_{0,4}$ 上，我们能找到这样的奇异黎曼曲面：它由两条 $\PP^1$ 粘在一起形成一个结点、且每条各带两个标记点。它们来自一条带四个标记点的光滑有理曲线：只需收缩一条把标记点分成二加二的圈。我们有三种可能的构型，对应于三种把 $(p_1,p_2,p_3,p_4)$ 分成两个互不相交、各含两点的子集的方式（见图 8）。

\figomit{7}{紧化模空间 $\Mc_{g,n}$ 的图示。}

% [原文 p.11]
注意到所有这些稳定黎曼曲面都没有模：它们正规化的每个有理分量都有三个特殊点（两个标记点与结点的一个分支），而这样三个点总可以被搬到 $(0,1,\infty)$。换一种说法：我们可以把上述每个稳定黎曼曲面都实现为点 $\cM_{0,3}\times\cM_{0,3}$。回忆 $\cM_{0,4}=\PP^1\setminus\{0,1,\infty\}$，我们得到
\begin{equation}
  \Mc_{0,4}=\cM_{0,4}\sqcup(\cM_{0,3}\times\cM_{0,3})^{\sqcup 3}=\PP^1,
  \tag{2.20}
\end{equation}
它确实是紧的。

% [Tutor Note: (2.20) 中指数上的 $\sqcup 3$ 表示“三个不交拷贝”。核对：把四点标号分成两个二元子集，不计两块顺序，共 $\frac12\binom{4}{2}=3$ 种分法，与图 8 的三条边界点一一对应，故边界部分是 $(\cM_{0,3}\times\cM_{0,3})^{\sqcup 3}$（每个因子 $\cM_{0,3}$ 是一点，故整个边界就是 3 个点）。而 $\cM_{0,4}=\PP^1\setminus\{0,1,\infty\}$ 本身就缺 3 个点，补上这 3 点恰好得到 $\PP^1$——这就是“它确实是紧的”的含义。]

至于 $\Mc_{1,1}$，边界 $\partial\Mc_{1,1}$ 中唯一的元素就是前面遇到过的那个带一个标记点的捏缩环面。同样，捏缩环面没有模，因为它的正规化是带三个标记点的有理曲线。然而，捏缩环面以 $\Z_2$ 为自同构群。换一种说法，就是把它实现为 $\cM_{0,3}/\Z_2$。这给出
\begin{equation}
  \Mc_{1,1}=\cM_{1,1}\sqcup(\cM_{0,3}/\Z_2),
  \tag{2.21}
\end{equation}
它在拓扑上是一个 $\PP^1$，但带有轨形结构：一个自同构为 $\Z_6$ 的点、一个自同构为 $\Z_4$ 的点，以及其余所有自同构为 $\Z_2$ 的点。

% [Tutor Note: 这里的三个轨形点可与 $j$-不变量对上号：一般点（$j\neq 0,1728$）的稳定子是 $\Z_2$（来自椭圆曲线的 $-1$），$j=1728$ 的点（即 $i$）稳定子是 $\Z_4$，$j=0$ 的点（即 $\rho=e^{2\pi\ii/3}$）稳定子是 $\Z_6$。原文按 $\Z_6,\Z_4,\Z_2$ 的顺序罗列并无特殊含义，但若对照 $\mathrm{PSL}(2,\Z)$ 的椭圆点阶数（$2,3$ 型）再看 $\mathrm{SL}(2,\Z)$ 的稳定子（阶 $2,4,6$），可确认点与群的对应关系不能互换。]

从上面的例子应当可以清楚看出，$\Mc_{g,n}$ 的紧化具有一种递归结构，它是通过收缩圈、并把黎曼曲面的拓扑拆解成若干块而得到的。

\figomit{8}{边界 $\partial\Mc_{0,4}$ 上的三个点。}

% [原文 p.12]
\figomit{9}{一个稳定黎曼曲面以及与之相伴的稳定图。}

我们可以借助某些图来追踪这一点。见图 9 的图示。

% [Tutor Note: 原文（以及数学文献）常把这里定义的稳定图称为“对偶图”（dual graph），见 p.13 “stable Riemann surface dual to the graph”。本译注本按 glossary 统一译作“稳定图”，只在需要强调“曲线与图互为对偶”时保留“对偶”二字。]

\begin{definition}[2.5]
  与稳定黎曼曲面 $(\Sigma,p_1,\dots,p_n)\in\Mc_{g,n}$ 相伴的\emph{稳定图}，是按下述方式得到的图 $\Gamma$：
  \begin{itemize}
    \item 给正规化的每个连通分量配一个顶点 $v$，并用该分量的亏格 $g(v)$ 来装饰它；
    \item 给每个标记点 $p_i$ 配一片叶子，并相应地标上标号 $i$；
    \item 给每个结点配一条边。
  \end{itemize}
\end{definition}

稳定图 $\Gamma$ 的亏格为
\begin{equation}
  g(\Gamma)=\sum_{v\in V(\Gamma)}g(v)+h^1(\Gamma),
  \tag{2.22}
\end{equation}
其中 $V(\Gamma)$ 是图的顶点集，$h^1(\Gamma)$ 表示 $\Gamma$ 的第一 Betti 数（即面的个数）。它与 $\Sigma$ 的亏格相等。我们还用 $E(\Gamma)$ 表示边集，用 $n(v)$ 表示顶点 $v$ 的度数（即与 $v$ 相接的叶子与半边的个数）。后者对应于与顶点 $v$ 相应的分量上特殊点（即标记点与结点的分支）的个数。

% [Tutor Note: 对连通图，$h^1(\Gamma)=|E(\Gamma)|-|V(\Gamma)|+1$，所以 (2.22) 可写成 $g(\Gamma)=\sum_v g(v)+|E|-|V|+1$。原文把 $h^1$ 释作“面的个数”，严格说这只在把图平面嵌入、且只数有界面时才对（Euler 公式 $V-E+F=2$ 给出有界面数 $F-1=E-V+1=h^{1}$）；一般图应读作“独立圈的个数”。核对方法：单顶点带一个自环（loop）：$V=1,E=1,h^{1}=1$，由 (2.22) 得 $g=g(v)+1$，正是把非分离结点光滑化后亏格减一。]
% [Tutor Note: 记号提醒：同一个字母 $n$ 在这里身兼二职——$(g,n)$ 中的 $n$ 是标记点总数，而 $n(v)$ 是顶点 $v$ 的度数。二者由 $\sum_{v}n(v)=n+2|E(\Gamma)|$ 相联系（每条边贡献两个半边），这个恒等式在习题 2.4 的维数计算中要用到。]

我们指出，稳定性条件蕴涵 $2g(v)-2+n(v)>0$ 对所有 $v\in V(\Gamma)$ 成立。这保证了对每个 $(g,n)$（称为\emph{类型}），只有有限多个亏格为 $g$、带 $n$ 片叶子的稳定图。这样的稳定图给出 $\Mc_{g,n}$ 的一个\emph{分层}：对给定类型 $(g,n)$ 的 $\Gamma$，置
\begin{equation}
  \cM_\Gamma=\left\{(\Sigma,p_1,\dots,p_n)\in\Mc_{g,n}\ \middle|\
  \begin{array}{l}
    \Gamma\ \text{是}\ (\Sigma,p_1,\dots,p_n)\\
    \text{所对应的稳定图}
  \end{array}
  \right\}.
  \tag{2.23}
\end{equation}
于是我们得到分层
\begin{equation}
  \Mc_{g,n}=\bigsqcup_{\Gamma\ \text{为}\ (g,n)\ \text{型}}\cM_\Gamma.
  \tag{2.24}
\end{equation}

% [Tutor Note: (2.24) 的并是“不交并”：一条稳定曲面只有一个稳定图，故不同的 $\cM_\Gamma$ 互不相交。核对退化情形：取 $|E(\Gamma)|=0$（没有结点），此时 $\Gamma$ 的顶点只有一个、亏格 $g$、$n$ 片叶子，$\cM_\Gamma$ 就是光滑层 $\cM_{g,n}$，它开稠密；其余 $\Gamma$ 给出边界分层。也正因如此，“分层”在这里不是代数几何中按维数递减的通常意义下的分层，而是把 $\Mc_{g,n}$ 分解成有限多个局部闭的子轨形。]

% [原文 p.13]
\figomit{10}{所有 $(0,4)$ 型与 $(1,1)$ 型的稳定图。}

我们已详细分析过 $\Mc_{0,4}$ 与 $\Mc_{1,1}$ 的情形，它们的稳定图在图 10 中给出。另一个例子是 $\Mc_2$，见图 11。那里我们画出了与图对偶的稳定黎曼曲面类型所对应的各分层，并按边数把它们画在不同的层级上。注意，圈的收缩对偶于边的收缩。

% [Tutor Note: “圈的收缩对偶于边的收缩”是 (2.22) 的几何对偶表述：曲面一侧收缩一个非分离圈使亏格减一，图一侧就收缩一条非分离边使 $h^{1}$ 减一、并把两个半边并成一个顶点，两边亏格同步下降。核对：在图 11 中，从上到下每降一层恰好多一条边，维数 $3,2,1,0$ 每次减一，与习题 2.4(2) 的公式 $\dim=\dim(\Mc_{g,n})-|E_\Gamma|$ 一致（$g=2$ 时 $\dim(\Mc_2)=3$）。]

\begin{exercise}[2.4]
  \begin{enumerate}
    \item 列出 $\Mc_{2,1}$ 的所有分层。
    \item 考虑类型为 $(g,n)$ 的稳定图 $\Gamma$。证明该分层的维数由下式给出：
    \[
      \dim(\cM_\Gamma)=\dim(\Mc_{g,n})-|E_\Gamma|.
    \]
  \end{enumerate}
\end{exercise}

% [Tutor Note: 习题 2.4(2) 的提示（勿写进正文）：用 (2.25) 的乘积公式算维数，$\dim\Mc_\Gamma=\sum_v(3g(v)-3+n(v))$，再用 $\sum_v n(v)=n+2|E_\Gamma|$ 与 $g=\sum_v g(v)+|E_\Gamma|-|V(\Gamma)|+1$（连通性）即得 $3g-3+n-|E_\Gamma|$。另注意原文此处把边集记作 $E_\Gamma$，而 p.12 定义处记作 $E(\Gamma)$——同一对象的两种写法，正文照排。]

$\Mc_{g,n}$ 的各分层由维数更小的空间 $\cM_\Gamma$ 参数化，这一事实有时被称为 $\Mc_{g,n}$ 的\emph{递归边界结构}。它是黎曼曲面模空间最重要的特征之一；关于 $\Mc_{g,n}$ 的许多结果（包括积分的计算）的证明都以一种非常本质的方式用到它。

\figomit{11}{亏格 2 的所有可能的稳定黎曼曲面，以及与之相伴的稳定图与相应分层的维数。}

% [原文 p.14]
利用它的一种方式是定义\emph{粘合映射}。更确切地说，对每个类型为 $(g,n)$ 的稳定图 $\Gamma$，我们定义
\begin{equation}
  \xi_\Gamma\colon\Mc_\Gamma=\prod_{v\in V(\Gamma)}\Mc_{g(v),n(v)}\longrightarrow\Mc_{g,n},
  \tag{2.25}
\end{equation}
它把稳定黎曼曲面 $((\Sigma_v)_{v\in V(\Gamma)},(q_h,q_{h'})_{e=(h,h')\in E(\Gamma)},p_1,\dots,p_n)$ 送到稳定黎曼曲面 $(\Sigma,p_1,\dots,p_n)$：后者由把所有对应于构成 $\Gamma$ 的边 $e$ 的半边对 $e=(h,h')$ 的点对 $(q_h,q_{h'})$ 粘合起来而得到。$\Mc_\Gamma$ 在 $\xi_\Gamma$ 下的像与 $\cM_\Gamma$ 的闭包重合。

% [Tutor Note: 记号辨析：(2.23) 的 $\cM_\Gamma$（花体 M，带下标 $\Gamma$）是一个分层，开而不闭；而 (2.25) 左边的 $\Mc_\Gamma$（带横线的 $\mathcal M$，带下标 $\Gamma$）是“乘积” $\prod_v\Mc_{g(v),n(v)}$，两者不是同一个空间。原文说“像与 $\cM_\Gamma$ 的闭包重合”，正是把二者的关系说清楚：$\xi_\Gamma(\Mc_\Gamma)=\overline{\cM_\Gamma}$。核对：$g=0,n=4$、$\Gamma$ 为一条分离边时，$\Mc_\Gamma=\Mc_{0,3}\times\Mc_{0,3}$ 是一点，$\xi_\Gamma$ 的像就是一个点（p.11 的三个边界点之一），而 $\overline{\cM_\Gamma}$ 也是这个点，因为 $|E_\Gamma|=1$、分层维数为 $0$。]

最简单的情形是恰含一条边 $e$ 的稳定图 $\Gamma$。有两种可能的情形：这条边是非分离的（即一个自环），或者它是。

% [Tutor Note: 原文此处句子残缺（“or it is.” 后无宾语），应为 “or it is separating (i.e. not a loop)”，即“或者它是分离的（即不是自环）”。判断依据：紧接着正文分两小节正是“Non-separating edge”与“Separating edge”两种情形，句子显然要引出这一二分。依“照排笔误”之例，正文保留残缺，正确写法在此与 ERRATA 中给出。]

\textbf{非分离边。}它对应于如下的稳定图：
\begin{equation}
  \text{〔稳定图：亏格 }g-1\text{ 的顶点，带 }n\text{ 个标记点 }1,\dots,n\text{ 与一条非分离边〕}
  \tag{2.26}
\end{equation}
因此，该粘合映射（称为\emph{非分离型}粘合映射）由下式给出：
\begin{equation}
  \begin{gathered}
    \rho\colon\Mc_{g-1,n+2}\longrightarrow\Mc_{g,n},\qquad\text{例如：}\\
    \begin{array}{c}
      \text{〔图：带标记点 }p_1\text{ 与两个待粘合的分支点 }q,q'\\
      \text{的曲面，把 }q,q'\text{ 粘合后得到带一个结点的曲面〕}
    \end{array}
  \end{gathered}
  \tag{2.27}
\end{equation}
严格说来，$\rho$ 应当依赖于 $(g,n)$。为了记号更轻便，我们省略这种依赖。

% [Tutor Note: 原文 (2.26)(2.27)(2.28)(2.29)(2.31) 的公式主体是插图，本译注本不复刻插图，改为在公式内用文字转述图上内容，编号与位置保持不变。核对 $\rho$ 的指标：源空间 $\Mc_{g-1,n+2}$ 的稳定性条件为 $2(g-1)-2+(n+2)=2g-2+n>0$，与目标一致；维数 $3(g-1)-3+(n+2)=3g-3+n-2$，比目标低 2（复一维），正是一个结点方向的模。]

\textbf{分离边。}它对应于如下的稳定图：
\begin{equation}
  \begin{array}{c}
    \text{〔稳定图：两个分别标有亏格 }g_1,g_2\text{ 的顶点，}\\
    \text{以一条分离边相连；左端带标记点集 }I_1\text{，右端带 }I_2\text{〕}
  \end{array}
  \tag{2.28}
\end{equation}
其中 $g=g_1+g_2$ 是亏格的一个分割，$I_1\sqcup I_2=\{p_1,\dots,p_n\}$ 是标记点的一个分割。于是相应的粘合映射（称为\emph{分离型}粘合映射）由下式给出：
\begin{equation}
  \begin{gathered}
    \sigma\colon\Mc_{g_1,1+|I_1|}\times\Mc_{g_2,1+|I_2|}\longrightarrow\Mc_{g,n},\qquad\text{例如：}\\
    \begin{array}{c}
      \text{〔图：带标记点 }p_1\text{ 与分支点 }q\text{ 的曲面，连同带分支点 }q'\\
      \text{与标记点 }p_2,p_3\text{ 的曲面；把 }q,q'\text{ 粘合后得到}\\
      \text{一个带结点、标记点为 }p_1,p_2,p_3\text{ 的曲面〕}
    \end{array}
  \end{gathered}
  \tag{2.29}
\end{equation}
严格说来，$\sigma$ 应当依赖于 $(g,n)$ 以及亏格与标记点的分割。

% [Tutor Note: $\sigma$ 的维数核对：源空间两因子的复维数之和为 $(3g_1-3+1+|I_1|)+(3g_2-3+1+|I_2|)=3(g_1+g_2)-3+(|I_1|+|I_2|)-2=3g-3+n-2$，同样比 $\dim\Mc_{g,n}$ 低 2，与“光滑化一个结点损失一个复模”相符。另外 $1+|I_1|$ 中的 $1$ 来自那条半边（它给出分量上的一个特殊点），这正是 (2.22) 之后所说的 $n(v)$ 计数方式。]

注意，上面的各项对应于拓扑递归公式（见 V. Bouchard 的讲义 [Bou26]）中出现的各项。这并非巧合，我们将在第 3 节看到这一点。

% [原文 p.15]
我们以模空间之间另一个自然的映射来结束本节：遗忘映射。这就是忘掉最后一个标记点的映射：
\begin{equation}
  \pi\colon\Mc_{g,n+1}\longrightarrow\Mc_{g,n},
  \qquad
  (\Sigma,p_1,\dots,p_n,p_{n+1})\longmapsto(\Sigma,p_1,\dots,p_n)^{\mathrm{stab}}.
  \tag{2.30}
\end{equation}
同样，严格说来 $\pi$ 应当依赖于 $(g,n)$。为了记号更轻便，我们省略这种依赖。后缀 “stab” 表示 “稳定化”（stabilisation）。的确，忘掉一个标记点时，所得的黎曼曲面可能是不稳定的。带三个特殊点的有理分量上的标记点 $p_{n+1}$ 就属于这种情况。稳定化过程只是把这个分量收缩成一点。如果所得的黎曼曲面仍然不稳定，我们就继续收缩不稳定分量，直到找到稳定的结果。例如：
\begin{equation}
  \begin{array}{c}
    \text{〔图：亏格 }1\text{ 的分量上带标记点 }p_1\text{，与之相连的}\\
    \text{有理分量上带标记点 }p_2,p_3\text{；遗忘 }p_3\text{ 后该有理分量}\\
    \text{不稳定（只剩两个特殊点），稳定化把它收缩成一点，}\\
    \text{得到带标记点 }p_1,p_2\text{ 的亏格 }1\text{ 曲面〕}
  \end{array}
  \tag{2.31}
\end{equation}
粘合映射 $\rho,\sigma$ 与遗忘映射 $\pi$ 有时统称为\emph{重言映射}。我们很快就会看到，它们在黎曼曲面模空间的相交理论中起着关键作用。要点在于：由于紧化以及重言映射的引入，我们可以把模空间看作由映射连接起来的一族空间（而不是 “孤立” 的）。特别地，我们可以在上同调中谈论拉回与前推。

% [Tutor Note: 稳定化可以反复进行，是因为 2.3 的稳定性条件是逐分量的 $2g-2+n>0$；一个亏格 $0$、只有两个特殊点的分量不满足它，必须缩掉。核对本例：开始时 $g=1,n=3$（$p_1$ 在亏格 1 分量上，$p_2,p_3$ 在附着的有理分量上），遗忘 $p_3$ 后该有理分量只剩 $p_2$ 与结点两个特殊点，不稳定，缩成一点后 $p_2$ 变成亏格 1 曲面上的标记点，得到 $g=1,n=2$。另外注意 $\pi$ 只在 $2g-2+n>0$、即目标空间存在时才有定义。]

让我们回顾对任意光滑映射的这些操作
\begin{equation}
  \phi\colon M\longrightarrow N
  \tag{2.32}
\end{equation}
其中 $M$ 与 $N$ 是实维数分别为 $\dim_{\R}(M)=m$ 与 $\dim_{\R}(N)=n$ 的光滑实轨形。

\textbf{拉回。}拉回总是一个定义良好的反变操作，它对应于预复合。更确切地说，它是一个保持次数的映射
\begin{equation}
  \phi^*\colon H^k(N)\longrightarrow H^k(M).
  \tag{2.33}
\end{equation}
用微分形式的语言，把 $\phi$ 局部地写成 $(x_1,\dots,x_m)\mapsto(y_1(x),\dots,y_n(x))$，并设 $\eta$ 是 $N$ 上的 $k$-形式，局部表示为 $\eta=\eta^{\mu_1,\dots,\mu_k}(y)\,\dd y_{\mu_1}\wedge\cdots\wedge\dd y_{\mu_k}$（我们对重复指标采用 Einstein 求和记号）。则
\begin{equation}
  \phi^*\eta=\eta^{\mu_1,\dots,\mu_k}(y(x))\,\dd y_{\mu_1}(x)\wedge\cdots\wedge\dd y_{\mu_k}(x).
  \tag{2.34}
\end{equation}
拉回与加法及 cup 积都相容。

% [Tutor Note: (2.34) 与正文的指标位置是自洽的：分量 $\eta^{\mu_1\cdots\mu_k}$ 用上标，微分 $\dd y_{\mu_i}$ 用下标，故求和是“上对下”的收缩；不要误改成 $\eta_{\mu_1\cdots\mu_k}\dd y^{\mu_i}$ 或混搭。又因为拉回与 cup 积相容，(2.33) 其实是分次环同态；这一点在下文把相交数写成拉回-前推的复合时至关重要。]

\textbf{前推。}前推只对纤维紧的映射 $\phi$ 才有良好定义。此时，前推给出上同调中的一个共变操作，它对应于 “沿纤维积分” 的几何想法。更确切地说，若用 $r$ 表示 $\phi$ 的纤维的实维数，则
\begin{equation}
  \phi_*\colon H^k(M)\longrightarrow H^{k-r}(N).
  \tag{2.35}
\end{equation}

% [Tutor Note: 前推的“次数下降 $r$”与“纤维紧”是同一件事的两面：沿 $r$ 维纤维积分把 $k$-形式变成 $(k-r)$-形式；纤维不紧时积分可能发散，“良好定义”就无从谈起。核对：$r=0$（即 $\phi$ 是有限覆叠）时 $\phi_*:H^k(M)\to H^k(N)$ 保持次数，与拉回同次数；而在开嵌入（纤维为点）的情形 $\phi_*$ 就是（可能带重数的）前推。用 Poincaré 对偶看，$\phi_*$ 对偶于 $\phi^{*}$（同调上），这也解释了为什么一个叫“共变”一个叫“反变”。]
